% knowledge-product-pdf-v6-components
\usepackage[most]{tcolorbox}
\usepackage{seqsplit}
\newcommand{\KpMixedGap}{\hskip .25em\relax}

% Adjacent glyph backgrounds form one tint while permitting line breaks.
% Zero horizontal padding preserves the mono advance and the original payload.
\newcommand{\KpCodeGlyph}[1]{{\small\ttfamily\setlength{\fboxsep}{0pt}\colorbox{kpBrandSoft}{\rule[-.20em]{0pt}{1.12em}#1}}}
\newtcbox{\KpCodeFirstBox}{on line,colback=kpBrandSoft,boxrule=0pt,boxsep=0pt,
  arc=.35em,sharp corners=east,left=.38em,right=0pt,top=0pt,bottom=0pt,
  nobeforeafter,tcbox raise base,fontupper=\small\ttfamily}
\newtcbox{\KpCodeLastBox}{on line,colback=kpBrandSoft,boxrule=0pt,boxsep=0pt,
  arc=.35em,sharp corners=west,left=0pt,right=.38em,top=0pt,bottom=0pt,
  nobeforeafter,tcbox raise base,fontupper=\small\ttfamily}
\newcommand{\KpCodeFirst}[1]{\KpCodeFirstBox{\rule[-.20em]{0pt}{1.12em}#1}}
\newcommand{\KpCodeLast}[1]{\KpCodeLastBox{\rule[-.20em]{0pt}{1.12em}#1}}
\newlength{\KpInlineWidth}
\newcommand{\KpCodeSmart}[2]{{%
  \xeCJKsetup{CJKecglue = \hskip 0pt}%
  \small\ttfamily\settowidth{\KpInlineWidth}{#1}%
  \advance\KpInlineWidth by 1.22em%
  \leavevmode\allowbreak\hspace{.18em}%
  \ifdim\KpInlineWidth<\linewidth
    \KpInlineCode{\vphantom{Ag}#1}%
  \else
    #2%
  \fi
  \hspace{.18em}%
}}

% Reserve a reader block using the active fonts and available line width.
\newdimen\KpMeasuredHeight
\newdimen\KpMeasuredWidth
\newdimen\KpMeasuredLineWidth
\newdimen\KpReserveLimit
\newcommand{\KpReserveBegin}[1]{\begingroup\KpMeasuredHeight=#1\relax\KpReserveLimit=\textheight}
% A long list/table/code may continue; reserve its start, not the entire page.
\newcommand{\KpReserveBreakable}{\KpReserveLimit=.55\textheight}
\newcommand{\KpMeasureLine}[1]{%
  \settowidth{\KpMeasuredWidth}{#1}%
  \advance\KpMeasuredHeight by \baselineskip
  \loop\ifdim\KpMeasuredWidth>\KpMeasuredLineWidth
    \advance\KpMeasuredHeight by \baselineskip
    \advance\KpMeasuredWidth by -\KpMeasuredLineWidth
  \repeat
}
\newcommand{\KpMeasureCodeLine}[1]{{%
  \small\ttfamily\linespread{1.10}\selectfont
  \KpMeasuredLineWidth=\dimexpr\linewidth-24pt\relax
  \KpMeasureLine{#1}%
  \global\KpMeasuredHeight=\KpMeasuredHeight
}}
\newcommand{\KpMeasureReaderText}[1]{{%
  \normalfont\normalsize
  \KpMeasuredLineWidth=\dimexpr\linewidth-20pt\relax
  \KpMeasureLine{#1}%
  \advance\KpMeasuredHeight by 5pt
  \global\KpMeasuredHeight=\KpMeasuredHeight
}}
\newcommand{\KpReserveEnd}{%
  \ifdim\KpMeasuredHeight<\KpReserveLimit
    \edef\KpReserveAction{\endgroup\noexpand\Needspace{\the\KpMeasuredHeight}}%
  \else
    \def\KpReserveAction{\endgroup\Needspace{6\baselineskip}}%
  \fi
  \KpReserveAction
}
% Measure the actual display formula, including its vertical skips. This only
% reserves space; the original Math node is still rendered once in the body.
\newcommand{\KpMeasureReaderBlock}[1]{{%
  \setbox0=\vbox{\hsize=\linewidth\normalfont\normalsize #1\par}%
  \advance\KpMeasuredHeight by \ht0
  \advance\KpMeasuredHeight by \dp0
  \global\KpMeasuredHeight=\KpMeasuredHeight
}}

% Measure finite closing table/paragraph combinations with the actual fonts.
% KpReserveEnd retains longtable continuation when the combination exceeds a page.
\newdimen\KpMeasuredRowHeight
\newcommand{\KpTableRowBegin}{\KpMeasuredRowHeight=1.28\baselineskip}
\newcommand{\KpMeasureTableCell}[3]{%
  \setlength{\KpMeasuredLineWidth}{(\linewidth-#2\tabcolsep)*\real{#1}}%
  \setbox0=\vbox{\hsize=\KpMeasuredLineWidth\normalfont\normalsize\raggedright #3\par}%
  \dimen0=\dimexpr\ht0+\dp0+4pt\relax
  \ifdim\dimen0>\KpMeasuredRowHeight \KpMeasuredRowHeight=\dimen0\fi
}

% knowledge-product-pdf-v6-inline-code
% Pandoc renders Markdown code spans as \texttt{...}; keep that semantic trigger
% and give every reader-visible inline code span a stable visual treatment.
\newtcbox{\KpInlineCode}{
  on line,
  colback=kpBrandSoft,
  colframe=kpBrand!25!white,
  boxrule=0.3pt,
  arc=0.35em,
  boxsep=0pt,
  left=0.38em,
  right=0.38em,
  top=0.12em,
  bottom=0.12em,
  nobeforeafter,
  tcbox raise base,
  fontupper=\small\ttfamily
}
% Lua Code targets only Markdown spans. URLs, math and Verbatim keep their roles.
\newtcolorbox{KpLongCodeBox}{
  breakable,colback=kpBrandSoft,colframe=kpBrand!25!white,boxrule=0.45pt,
  arc=2pt,left=6pt,right=6pt,top=4pt,bottom=4pt,fontupper=\small\ttfamily\xeCJKsetup{CJKecglue = \hskip 0pt},
  before skip=5pt,after skip=5pt
}
\newcommand{\KpLongInlineCode}[1]{\seqsplit{#1}}
\AtBeginDocument{\pdfstringdefDisableCommands{\def\KpInlineCode#1{#1}\def\KpLongInlineCode#1{#1}\def\KpCodeSmart#1#2{#1}}}

\newtcolorbox{KpChapterLead}{
  enhanced,
  colback=kpBrandSoft,colframe=kpBrandSoft,boxrule=0pt,
  borderline west={2pt}{0pt}{kpBrand},arc=2pt,
  left=10pt,right=10pt,top=8pt,bottom=8pt,
  before skip=0pt,after skip=12pt,
  after={\Needspace{5\baselineskip}}
}

\newtcolorbox{KpSourceCard}{
  breakable,
  colback=kpPaper,
  colframe=\KnowledgeProductSourceFrameColor,
  boxrule=0pt,
  borderline west={0.6pt}{0pt}{kpRule},
  arc=2pt,
  left=10pt,
  right=10pt,
  top=3pt,
  bottom=3pt,
  fontupper=\small,
}
\newtcolorbox{KpKeyTakeaway}{
  breakable,
  colback=kpBrandSoft,
  colframe=\KnowledgeProductKeyFrameColor,
  boxrule=0.5pt,
  arc=2pt,
  left=10pt,
  right=10pt,
  top=8pt,
  bottom=8pt,
}
\newtcolorbox{KpPractice}{
  breakable,
  colback=kpPaper,
  colframe=kpRule,
  boxrule=0.5pt,
  arc=2pt,
  left=10pt,
  right=10pt,
  top=8pt,
  bottom=8pt,
}
\newtcolorbox{KpNote}{
  breakable,
  colback=kpPaper,
  colframe=kpRule,
  boxrule=0.5pt,
  arc=2pt,
  left=10pt,
  right=10pt,
  top=8pt,
  bottom=8pt,
}
\newtcolorbox{KpDecision}{
  breakable,
  colback=kpBrandSoft,
  colframe=\KnowledgeProductKeyFrameColor,
  boxrule=0.5pt,
  arc=2pt,
  left=10pt,
  right=10pt,
  top=8pt,
  bottom=8pt,
}
\newtcolorbox{KpProcedure}{
  breakable,
  colback=kpPaper,
  colframe=kpBrand,
  boxrule=0.5pt,
  arc=2pt,
  left=10pt,
  right=10pt,
  top=8pt,
  bottom=8pt,
}
\newtcolorbox{KpWorksheet}{
  breakable,
  colback=kpPaper,
  colframe=kpRule,
  boxrule=0.7pt,
  arc=2pt,
  left=10pt,
  right=10pt,
  top=10pt,
  bottom=10pt,
}
\newtcolorbox{KpPercentageDistribution}{
  breakable,
  colback=kpBrandSoft,
  colframe=kpRule,
  boxrule=0.5pt,
  arc=2pt,
  left=10pt,
  right=10pt,
  top=8pt,
  bottom=8pt,
}
\newtcolorbox{KpCodeBox}{
  breakable,
  title after break={代码（续）},fonttitle=\small\bfseries,
  coltitle=kpInk,colbacktitle=kpPaper,
  colback=kpPaper,
  colframe=kpRule,
  boxrule=0.5pt,
  arc=3pt,
  left=8pt,
  right=8pt,
  top=8pt,
  bottom=8pt,
}
